% Figure: strain energy density as the area under the stress-strain line.
% Linear elastic loading: u = 1/2 sigma epsilon. The triangle under the
% line is the recoverable elastic energy per unit volume; the resilience
% is the triangle up to yield.
\begin{figure}[htbp]
\centering
\begin{tikzpicture}
\begin{axis}[
  width=11cm, height=8cm,
  xlabel={strain $\epsilon$},
  ylabel={stress $\sigma$},
  xmin=0, xmax=1.15, ymin=0, ymax=1.15,
  axis lines=left,
  xtick=\empty, ytick=\empty,
  clip=false,
]
% elastic line up to yield
\addplot[engblue, very thick, domain=0:0.8, samples=2] {1.0*x};
% shaded resilience triangle up to yield
\addplot[englime!35, draw=none] coordinates {(0,0) (0.8,0.8) (0.8,0)} \closedcycle;
% generic operating point
\addplot[engred, very thick, domain=0:0.55, samples=2] {1.0*x};
\addplot[engred!30, draw=none] coordinates {(0,0) (0.55,0.55) (0.55,0)} \closedcycle;
% dashed projections at operating point
\draw[engred, dashed] (axis cs:0.55,0) -- (axis cs:0.55,0.55);
\draw[engred, dashed] (axis cs:0,0.55) -- (axis cs:0.55,0.55);
% yield marker
\draw[enggray, dashed] (axis cs:0.8,0) -- (axis cs:0.8,0.8);
\draw[enggray, dashed] (axis cs:0,0.8) -- (axis cs:0.8,0.8);
% labels
\node[engred, font=\scriptsize, anchor=west] at (axis cs:0.30,0.13)
  {$u=\tfrac{1}{2}\sigma\epsilon$};
\node[englime!50!black, font=\scriptsize, anchor=south east]
  at (axis cs:0.78,0.30) {resilience $u_R=\dfrac{\sigma_y^2}{2E}$};
\node[engred, font=\scriptsize, anchor=south] at (axis cs:0.55,0.57)
  {$(\epsilon,\sigma)$};
\node[enggray, font=\scriptsize, anchor=south] at (axis cs:0.8,0.82)
  {$(\epsilon_y,\sigma_y)$};
\node[engblue, font=\scriptsize, anchor=west] at (axis cs:0.82,0.70)
  {slope $E$};
\end{axis}
\end{tikzpicture}
\caption[Strain energy density as the area under the elastic line]{Strain
energy density is the area under the stress-strain curve. For linear
elastic loading the area to a generic point is the triangle
$u=\tfrac{1}{2}\sigma\epsilon=\tfrac{1}{2}E\epsilon^2=\sigma^2/(2E)$, the
recoverable energy stored per unit volume (red). The area to first yield
is the \emph{modulus of resilience} $u_R=\sigma_y^2/(2E)$ (green), the
energy a unit volume of material can store and return without permanent
set. A spring is engineered to make this triangle large: a high yield
stress and a moderate modulus store the most recoverable energy per unit
volume.}
\label{fig:vol03ch03:strain-energy-triangle}
\end{figure}
